\documentclass[10pt,journal,compsoc]{IEEEtran}

\usepackage{cite}
\usepackage{amsmath,amssymb,amsfonts}
\usepackage{algorithmic}
\usepackage{graphicx}
\usepackage{textcomp}
\usepackage{xcolor}
\usepackage{booktabs}
\usepackage{hyperref}

\begin{document}

\title{An Intelligent Multi-Modal Weather Intelligence and Climate Decision Support Platform for Indian Southwest Monsoon Analysis}

\author{MacroEdtech Team GenAI Research Program 2026\\
\textit{MacroEdtech Educational Initiative}\\
Indore, Madhya Pradesh, India\\
\texttt{info@macroedtech.net}
}

\maketitle

\begin{abstract}
The Southwest Monsoon is the critical driver of socio-economic and agricultural stability in India. Precise meteorological prediction and accessible climate insights are essential for disaster risk reduction and evidence-based policy formulation. This paper presents an open-source, multi-modal Weather Intelligence and Climate Decision Support Platform. The system combines statistical machine learning, deep learning time-series architectures, computer vision for remote sensing satellite analysis, and an Agentic Retrieval-Augmented Generation (RAG) framework. By integrating predictive engines with open-weight Large Language Models (LLMs) and Model Context Protocol (MCP) agents, the platform enables domain-grounded automated forecasting, literature retrieval, and decision support.
\end{abstract}

\begin{IEEEkeywords}
Generative AI, Southwest Monsoon, Retrieval-Augmented Generation (RAG), Time-Series Forecasting, Deep Learning, Remote Sensing, AI Agents.
\end{IEEEkeywords}

\section{Introduction}
The Indian Southwest Monsoon governs precipitation patterns across South Asia. Precise forecasting of extreme weather, rainfall spatial distribution, and temperature anomalies remains a significant challenge due to chaotic atmospheric dynamics. Traditional numerical weather prediction (NWP) models demand significant computational resources and often lack real-time interpretability for non-technical stakeholders.

This project introduces an end-to-end open-source AI platform designed to bridge predictive data science and actionable climate decision-making. The system integrates predictive machine learning and deep learning pipelines with modern Generative AI engineering, enabling natural language interactions, scientific report generation, and multi-agent workflow orchestration.

\section{System Architecture}
The overall platform architecture is organized into two primary subsystems:

\subsection{Part 01: Multi-Modal Weather Prediction Engine}
The prediction pipeline processes multi-source meteorological observations and satellite datasets, including ERA5 Reanalysis, Open-Meteo APIs, and Sentinel-2 satellite imagery.
\begin{itemize}
    \item \textbf{Machine Learning \& Time-Series:} Classical algorithms (XGBoost, LightGBM, Random Forest) alongside advanced time-series networks (LSTM, CNN-LSTM, Temporal Fusion Transformer).
    \item \textbf{Computer Vision:} Image segmentation and spatial feature extraction using U-Net and ResNet architectures for cloud-cover and remote sensing analytics.
    \item \textbf{Backend Deployment:} Microservices built using FastAPI, fully containerized via Docker.
\end{itemize}

\subsection{Part 02: Generative AI \& Agentic Decision Support}
To transform predictive outputs into readable climate intelligence, the platform employs a domain-specific GenAI architecture:
\begin{itemize}
    \item \textbf{Knowledge Base \& Vector DB:} Scientific literature and meteorological reports indexed via dense vector embeddings using ChromaDB.
    \item \textbf{RAG Pipeline:} Contextual grounding using open-weight LLMs (e.g., Llama 3, Gemma, Mistral) to minimize hallucinations.
    \item \textbf{Agentic Workflows:} Orchestrated with LangGraph and tool calling protocols, allowing AI agents to query predictive APIs dynamically.
\end{itemize}

\section{Experimental Setup and Results}
Models were evaluated across multiple standard statistical metrics, including Mean Absolute Error (MAE), Root Mean Squared Error (RMSE), and the Coefficient of Determination ($R^2$). 

The empirical evaluation results across tested model architectures are summarized in Table~\ref{tab:model_performance}.

% Dynamic Table Import generated by scripts/export_latex.py
\begin{table}
\caption{Comparative Performance Evaluation of Monsoon Prediction Models}
\label{tab:model_performance}
\begin{tabular}{|l|c|c|c|}
\toprule
MODEL ARCHITECTURE & MAE & RMSE & R2 \\
\midrule
XGBoost Baseline & 2.1400 & 3.8200 & 0.8420 \\
CNN-LSTM Hybrid & 1.4500 & 2.1100 & 0.9150 \\
Temporal Fusion Transformer & 1.1200 & 1.7800 & 0.9480 \\
\bottomrule
\end{tabular}
\end{table}


\section{Discussion and Limitations}
Experimental results indicate that hybrid deep learning models (such as CNN-LSTM and TFT) outperform isolated statistical baselines when capturing complex monsoon spatio-temporal dependencies. Integrating RAG ensures that climate recommendations generated by the LLM remain grounded in official research publications and verified data sources.

\textbf{Limitations:} Performance remains dependent on local dataset completeness and regional sensor density. High-resolution satellite image processing introduces latency during live inference.

\section{Conclusion and Future Work}
This work demonstrates a production-ready, open-source Weather Intelligence Platform integrating ML predictions with Generative AI agent orchestration. Future work will expand multi-modal vision capabilities for real-time flood monitoring and extend edge deployment optimizations.

\section*{Acknowledgment}
The author acknowledges the mentorship and guidance provided by Sagar Sakalley (Founder \& Director, MacroEdtech) during Phase 02 of the Team GenAI Research Program 2026.

\end{document}